%% ---------------------------------------------------------------
%% table 객체: tab:audit
%% 파일: tables/tab_audit.tex
%% 사용처: chapters/10_discussion/03_audit.tex  (%% ---------------------------------------------------------------
%% table 객체: tab:audit
%% 파일: tables/tab_audit.tex
%% 사용처: chapters/10_discussion/03_audit.tex  (%% ---------------------------------------------------------------
%% table 객체: tab:audit
%% 파일: tables/tab_audit.tex
%% 사용처: chapters/10_discussion/03_audit.tex  (%% ---------------------------------------------------------------
%% table 객체: tab:audit
%% 파일: tables/tab_audit.tex
%% 사용처: chapters/10_discussion/03_audit.tex  (\input{tables/tab_audit})
%% 규칙: 캡션·라벨·내용을 이 파일 안에서만 관리한다. 수치는 config/numbers.tex 매크로를 쓴다.
%% ---------------------------------------------------------------
\begin{table}[htbp]
  \centering
  \caption{사전 감사에서 수정된 코드--논문 불일치}
  \label{tab:audit}
  \small
  \begin{tabularx}{\linewidth}{@{}lXX@{}}
    \toprule
    항목 & 발견 & 수정 \\
    \midrule
    A3 지름 분할 & 단일 연결(4{,}000 단위)로 분할하여 사슬형 군집의 지름이 상한을 초과할 수 있었음 & 완전 연결·지름 상한 분할로 교체 \\
    A1 온셋 검증 & ``생존 2 명(지도 전체)''과 ``반경 내 2 명(생존 무관)''을 별개 조건으로 검사 & 생존 $\wedge$ 반경 내를 단일 논리곱으로 검사 \\
    A2 인접 병합 & 15 초 / 2{,}000 단위 병합이 설정만 있고 실행되지 않음 & 병합 단계 구현 \\
    B1 특수 킬 & \texttt{CHAMPION\_SPECIAL\_KILL}을 두 번째 킬로 계수 & 보너스만 반영 \\
    MLP 경로 & 헤드라인 러너의 MLP가 95 차원 거시 시퀀스를 사용 & 2{,}980 차원 정합 입력으로 재경로 \\
    \bottomrule
  \end{tabularx}
\end{table}
)
%% 규칙: 캡션·라벨·내용을 이 파일 안에서만 관리한다. 수치는 config/numbers.tex 매크로를 쓴다.
%% ---------------------------------------------------------------
\begin{table}[htbp]
  \centering
  \caption{사전 감사에서 수정된 코드--논문 불일치}
  \label{tab:audit}
  \small
  \begin{tabularx}{\linewidth}{@{}lXX@{}}
    \toprule
    항목 & 발견 & 수정 \\
    \midrule
    A3 지름 분할 & 단일 연결(4{,}000 단위)로 분할하여 사슬형 군집의 지름이 상한을 초과할 수 있었음 & 완전 연결·지름 상한 분할로 교체 \\
    A1 온셋 검증 & ``생존 2 명(지도 전체)''과 ``반경 내 2 명(생존 무관)''을 별개 조건으로 검사 & 생존 $\wedge$ 반경 내를 단일 논리곱으로 검사 \\
    A2 인접 병합 & 15 초 / 2{,}000 단위 병합이 설정만 있고 실행되지 않음 & 병합 단계 구현 \\
    B1 특수 킬 & \texttt{CHAMPION\_SPECIAL\_KILL}을 두 번째 킬로 계수 & 보너스만 반영 \\
    MLP 경로 & 헤드라인 러너의 MLP가 95 차원 거시 시퀀스를 사용 & 2{,}980 차원 정합 입력으로 재경로 \\
    \bottomrule
  \end{tabularx}
\end{table}
)
%% 규칙: 캡션·라벨·내용을 이 파일 안에서만 관리한다. 수치는 config/numbers.tex 매크로를 쓴다.
%% ---------------------------------------------------------------
\begin{table}[htbp]
  \centering
  \caption{사전 감사에서 수정된 코드--논문 불일치}
  \label{tab:audit}
  \small
  \begin{tabularx}{\linewidth}{@{}lXX@{}}
    \toprule
    항목 & 발견 & 수정 \\
    \midrule
    A3 지름 분할 & 단일 연결(4{,}000 단위)로 분할하여 사슬형 군집의 지름이 상한을 초과할 수 있었음 & 완전 연결·지름 상한 분할로 교체 \\
    A1 온셋 검증 & ``생존 2 명(지도 전체)''과 ``반경 내 2 명(생존 무관)''을 별개 조건으로 검사 & 생존 $\wedge$ 반경 내를 단일 논리곱으로 검사 \\
    A2 인접 병합 & 15 초 / 2{,}000 단위 병합이 설정만 있고 실행되지 않음 & 병합 단계 구현 \\
    B1 특수 킬 & \texttt{CHAMPION\_SPECIAL\_KILL}을 두 번째 킬로 계수 & 보너스만 반영 \\
    MLP 경로 & 헤드라인 러너의 MLP가 95 차원 거시 시퀀스를 사용 & 2{,}980 차원 정합 입력으로 재경로 \\
    \bottomrule
  \end{tabularx}
\end{table}
)
%% 규칙: 캡션·라벨·내용을 이 파일 안에서만 관리한다. 수치는 config/numbers.tex 매크로를 쓴다.
%% ---------------------------------------------------------------
\begin{table}[htbp]
  \centering
  \caption{사전 감사에서 수정된 코드--논문 불일치}
  \label{tab:audit}
  \small
  \begin{tabularx}{\linewidth}{@{}lXX@{}}
    \toprule
    항목 & 발견 & 수정 \\
    \midrule
    A3 지름 분할 & 단일 연결(4{,}000 단위)로 분할하여 사슬형 군집의 지름이 상한을 초과할 수 있었음 & 완전 연결·지름 상한 분할로 교체 \\
    A1 온셋 검증 & ``생존 2 명(지도 전체)''과 ``반경 내 2 명(생존 무관)''을 별개 조건으로 검사 & 생존 $\wedge$ 반경 내를 단일 논리곱으로 검사 \\
    A2 인접 병합 & 15 초 / 2{,}000 단위 병합이 설정만 있고 실행되지 않음 & 병합 단계 구현 \\
    B1 특수 킬 & \texttt{CHAMPION\_SPECIAL\_KILL}을 두 번째 킬로 계수 & 보너스만 반영 \\
    MLP 경로 & 헤드라인 러너의 MLP가 95 차원 거시 시퀀스를 사용 & 2{,}980 차원 정합 입력으로 재경로 \\
    \bottomrule
  \end{tabularx}
\end{table}
